% 本文件由 LaTeX 排版 Agent 根据有效 translation.json 自动生成。
% author=Athanasius; work_id=2819; title=致非洲人会议书信

\chapter{\WorkTitle{致非洲人会议书信}\OriginalTerm{}{Ad Afros Epistola Synodica}}

% structure: p0001 source=h1 target=omitted reason=page-title-already-rendered-by-parent

\Emph{致非洲主教们}\par

\Emph{九十位埃及及利比亚主教（包括\Person{亚大纳削}）的书信}\par

1. \Emph{尼西亚大公会议的卓越地位。试图抬举\Place{里米尼}会议而贬低它的努力。}\par

我们所爱的同道、\Place{罗马}的主教\Person{达玛苏}，以及众多与他一同聚会的同道，他们的信件已足够了；同样，在\Place{高卢}和\Place{意大利}举行的其他会议，论及基督所赐、宗徒所宣讲、由来自全世界聚集在\Place{尼西亚}的教父们所传承的纯正\OriginalTerm{信德}{faith}，它们的信件也一样足够。当时因\OriginalTerm{阿里乌斯异端}{Arian heresy}引起极大的骚动，为使陷入其中者得以回头，并使\OriginalTerm{异端}{heresy}的始作俑者显露出来。因此，整个世界早已认同那次会议。如今，尽管举行过许多会议，在\Place{达尔马提亚}和\Place{达尔达尼亚}、\Place{马其顿}、\Place{伊庇鲁斯}和\Place{希腊}、\Place{克里特}及其他岛屿、\Place{西西里}、\Place{塞浦路斯}、\Place{潘菲利亚}、\Place{吕基亚}和\Place{伊索里亚}、整个\Place{埃及}和\Place{利比亚}，以及大部分\Place{阿拉伯}地区，所有人都已知道了它，并对那些签署者表示惊奇；即使他们中间仍存有那从\OriginalTerm{阿里乌斯派}{Arians}根上萌发的苦毒——我们指的是\Person{奥克森提乌斯}、\Person{乌尔撒基乌斯}、\Person{瓦伦斯}及其同党——他们也因这些信件而被断绝和孤立。我们再说一次：在\Place{尼西亚}达成的\OriginalTerm{宣认}{confession}，本身已足够且充足，足以颠覆一切不虔敬的异端，并维护和发展教会的道理。但我们听说，某些意欲反对它的人，企图援引一次据称在\Place{里米尼}举行的会议，并热切地要使它胜于那次会议；因此，我们写信提醒你们，绝不可容忍此类事：因为这不外是阿里乌斯异端的再次萌发。那些拒绝反抗异端的会议——即尼西亚会议——的人，若不是想让阿里乌斯的事业得胜，他们又想望什么呢？那么，这样的人除了被称为阿里乌斯派，并与阿里乌斯派同受惩罚外，还配得什么？因为祂说过：\BibleQuote{不可移动你祖先所立的古界}（见：\BibleRef{箴 22:28}），及\BibleQuote{凡咒骂父亲或母亲的，应处死刑}（见：\BibleRef{出 21:17}），他们却不敬畏天主；他们也不敬畏那些训令将所有持与之信德宣认相反意见者\OriginalTerm{绝罚}{anathema}的教父们。\par

2. \Emph{\Place{尼西亚}会议与此后历次地方会议的比较}\par

正因如此，在\Place{尼西亚}举行了\OriginalTerm{大公会议}{ecumenical synod}，318位主教因阿里乌斯异端聚集讨论\OriginalTerm{信德}{faith}；其目的就在于：关于信德，不应再举行地方会议，即使举行，也应视为无效。因为那次大公会议有何不足，竟需有人妄想革新？它充满了虔诚，亲爱的诸位，且已使全世界充满虔诚；\Place{印度}人已承认它，其他所有蛮族的基督徒也都承认了。因此，那些多次试图反抗它的人，他们的劳苦都是徒然的。因为我们所指的那些人已经召开了十次或更多会议，每次改变立场，从先前的决议中删减一些，在后来者中又做更改或增加。至今他们已藉着写作、删除和施加武力而一无所获，却不知道：\BibleQuote{凡是天父所没有种植的植物，必被拔除}（见：\BibleRef{玛 15:13}）。\BibleQuote{但透过尼西亚大公会议所传达的主的话，永远常存}（见：\BibleRef{伯前 1:25}）。若将人数与人数相比，在尼西亚聚会者多于地方会议，因为全体大于部分。但若有人想分辨尼西亚会议的原因，以及这些人后来召开的众多会议的原因，他就会发现：前者有合理的原因，后者却是因仇恨和争竞被武力拼凑而成的。因为前次会议是因阿里乌斯异端以及复活节而召开：当时\Place{叙利亚}、\Place{奇里乞亚}和\Place{美索不达米亚}的人与我们不同，他们在与犹太人相同的时节守此节。但感谢主，不仅在信德上，而且在\OriginalTerm{圣节}{Sacred Feast}上，已达成了和谐。这就是尼西亚会议的原因。但后来的会议不可胜数，全都是针对大公会议而策划的。\par

3. \Emph{\Place{里米尼}会议的真实性质}\par

这一点被指出之后，谁还会接受那些引用\Place{里米尼}会议或其他会议来反对\Place{尼西亚}会议的人呢？谁又能不憎恨那些轻视教父们的决议，而将那些在\Place{里米尼}以纷争与暴力制定的新决议置于其上的人呢？谁愿意赞同这些连自己的决议都不接受的人呢？因为正如我前面说过的，在他们自己召开的十次或更多次会议中，他们时而写这个，时而写那个，结果显然成了互相控告的人。他们的情形与古时犹太叛徒的情形相似。正如他们离弃了活水的泉源，为自己凿了破裂不能存水的池子，正如先知 \Person{耶肋米亚} 所说的那样，同样这些人，攻击那个大公会议，\BibleQuote{为自己凿出}许多会议，而所有这些会议都显为空虚，如同\Quote{没有力量的禾捆}。因此，我们不要容忍那些引用\Place{里米尼}或其他会议来反对\Place{尼西亚}会议的人。因为那些引用\Place{里米尼}会议的人，似乎不知道在那里发生的事，否则他们就不会提它了。因为，亲爱的弟兄们，你们知道，从你们那里前往\Place{里米尼}的人，知道\Person{乌尔撒奇乌斯}、\Person{瓦伦斯}、\Person{欧多克修斯}和\Person{奥克森提乌斯}（当时\Person{德默斐洛}也与他们一起），在企图写下某些取代\Place{尼西亚}决议的东西之后，是如何被罢免的。因为当他们被要求谴责亚略异端时，他们拒绝了，宁愿做这异端的首领。因此，主教们，作为主的真正仆人和正统的信徒（当时将近有二百位），写明他们只满足于\Place{尼西亚}会议，并希望和持守的既不多也不少，正是这一点。他们也将此报告给了皇帝\Person{君士坦提乌斯}，正是他下令召开会议的。然而，那些在\Place{里米尼}被罢免的人，去到\Person{君士坦提乌斯}那里，使那些反对他们的人受到侮辱，并受到威胁，不许返回各自的教区，且于同年冬天在\Place{色雷斯}遭暴力对待，以强迫他们容忍这些人的创新。\par

4. \Emph{符合圣经的尼西亚信式}\par

因此，如果有人引用\Place{里米尼}会议，首先请他们指出上述人士被罢免的事实，以及主教们所写的内容，即任何人都不应寻求超出尼西亚教父们所同意的内容，也不应引用那会议以外的任何会议。然而他们却隐瞒这些，反而大肆渲染在\Place{色雷斯}以暴力所行之事，从而表明自己既是亚略异端的伪善者，又是健全信仰的背离者。再者，如果一个人去审查并比较那大公会议本身与这些人所召开的会议，他就会发现前者的虔诚与后者的愚昧。在\Place{尼西亚}聚集的人们并非在被罢免之后才开会；其次，他们信奉子是出于圣父的\OriginalTerm{性体}{Essence}。但是另一些人，在屡次被罢免之后，并在\Place{里米尼}会议上再次被罢免，竟胆敢写下不应说子具有性体或\OriginalTerm{本体}{Subsistence}。弟兄们，这使我们看清，\Place{尼西亚}会议的人们散发出圣经的精神，因为天主在\BibleBook{出谷}{纪}中说：\BibleQuote{我是自有者}（\BibleRef{出 3:14}），又藉\Person{耶肋米亚}说：\BibleQuote{谁曾站在他的实体中，看见他的话呢？}紧接着又说：\BibleQuote{如果他们站在我的本位中，听我的话。}现在，\OriginalTerm{本体}{Subsistence}就是性体，无非就是存在本身，\Person{耶肋米亚}称之为\Quote{存在}，正如他说：\BibleQuote{他们不听存在的声音。}因为本体和性体就是存在：因为它存在，或换言之，它是。\Person{保禄}也领悟了这一点，在\BibleBook{致希伯来人}{书}中写道：\BibleQuote{他是天主光荣的辉耀，是天主本体的真像}（\BibleRef{希 1:3}）。但是那些自以为认识圣经并自称智者、不愿提及在天主内的本体的人（因为他们曾在\Place{里米尼}及其它会议上如此写下），当他们如同愚人在心中说：\Quote{没有天主}时，确实被公义地罢免了。而且，\Place{尼西亚}的教父们教导，子和圣言并非受造物，也非被造，因为他们读过：\BibleQuote{万物是藉着他而造成的}（\BibleRef{若 1:3}），以及\BibleQuote{一切都是在他内受造的，并在他内而存在}（\BibleRef{哥 1:16}）。而这些人，与其说是基督徒，不如说是亚略派，在他们的其它会议上，竟敢称呼他为受造物，以及受造之物中的一员，而他自己正是这些万物的工匠和创造者。因为如果\BibleQuote{万物是藉着他而造成的}，而他本身也是受造物，那么他将成为自己的创造者。然而，那正在受造的怎能创造？或那正在创造的岂能受造？\par

5. \Emph{尼西亚如何采纳\OriginalTerm{同性同体}{Coessential}的检验}\par

但他们即使这样说，也不感到羞耻，尽管他们说的话使所有人都憎恨他们；他们援引\Place{阿里米努姆}会议，只是为了表明他们也在那里被罢免了。至于\Place{尼西亚}的实际定义，即子与父\OriginalTerm{同性同体}{coessential}，他们表面上为此反对会议，并像蚊蚋一样围绕着这个表述到处嗡嗡作响，他们要么出于无知而绊倒，就如那些在\Place{熙雍}放置的绊脚石上绊倒的人（\BibleRef{罗 9:33}）；要么他们明明知道，却正是为此坚持不懈地反对和抱怨，因为这定义是准确的宣言，并且直面他们的异端。因为令他们恼怒的不是那些用语，而是定义中所包含的对他们自身的谴责。而这，再次，是他们自己造成的，即使他们想隐瞒他们完全清楚的事实——但我们现在必须提及这一点，以便由此也可展示伟大公会议的准确性。因为聚集的主教们想要摒弃亚略派编造的亵渎用语，即\Quote{从无中造成}，以及子是\Quote{受造之物}、\Quote{受造物}，以及\Quote{曾有一段时间他并不存在}，以及\Quote{他具有可变的本性}。并且他们想书面记录圣经公认的语言，即圣言按本性是天主的独生子、德能、智慧、真天主，如\Person{若望}所说，\Person{保禄}所写，他是父光荣的辉耀和他本体的真像。但\Person{优西比乌}及其同伙，被自己的错谬牵引，持续商议如下：\Quote{让我们同意吧。因为我们也是出于天主：因为\BibleQuote{只有一个天主，万物都出于他}\BibleRef{格前 8:6}，和\BibleQuote{旧的已成过去，看，都成了新的，但一切都是出于天主}\BibleRef{格后 5:17}。}他们又思考\WorkTitle{牧者}中所写的：\Quote{在一切之前，要相信天主是唯一的，他创造并安排了万物，并使它们从无中而存在。}但主教们，看穿了他们的狡诈和他们不虔的狡猾，更清晰地表达了\Quote{出于天主}这话的意义，他们写明子是\OriginalTerm{出自天主的本质}{of the Essence of God}，这样，受造物，由于不是无原因地自己存在，而是有存在的开端，被说成是\Quote{出于天主}，而唯独子可被视为专属父的本质。因为这是唯一的子和真实的圣言与父的关系所独有的，这也是采用\Quote{\OriginalTerm{出自本质}{of the essence}}这一用语的理由。再者，当主教们问那些伪装的少数派是否同意子不是受造物，而是父的德能和唯一智慧，是父的永恒且各方面完全的真像和真天主时，人们看到\Person{优西比乌}及其同伙互相点头示意，仿佛在说\Quote{这也适用于我们众人，因为我们也被称为\BibleQuote{天主的肖像和光荣}\BibleRef{格前 11:7}，并且论及我们说，\InnerQuote{我们这些活着的人，时常，}还有许多大能者，而\InnerQuote{上主的一切大能都出了埃及地，}毛虫和蝗虫被称为他的\BibleQuote{大能}\BibleRef{岳 2:25}。并且\InnerQuote{万军的上主与我们同在，雅各伯的天主是我们的护佑。}因为我们坚持我们是属于天主的，而且不仅如此，他甚至称我们为兄弟。即使他们称子为真天主，我们也不烦恼。因为他被造时，就实在地存在了。}\par

\Emph{尼西亚的检验在意义上并非不符合圣经，亦非新奇之事。}\par

亚略派的心智如此败坏。但在这里，主教们也看穿了他们的狡诈，从圣经中收集了光辉、河与井的比喻，以及\OriginalTerm{真像}{express Image}与\OriginalTerm{本体}{Subsistence}关系的经文，还有\BibleQuote{在你的光中，我们必得见光}和\BibleQuote{我与父原是一体}\BibleRef{若 10:30}等经文。最后他们更清晰简洁地写明，子与父\OriginalTerm{同性同体}{coessential}；因为上述所有经文都意味着这一点。而他们的抱怨，说这些用语不符合圣经，被他们自己证明是徒劳的，因为他们已经用不符合圣经的用语表达了他们的不虔：（例如\Quote{从无中}和\Quote{曾有一段时间他不存在}），然而他们却挑剔，因为他们被意义虔敬却不符合圣经的用语定罪。他们就像从粪堆里出来的人，实实在在是\BibleQuote{讲论地的事}\BibleRef{若 3:31}，而主教们并非自己发明这些用语，而是从教父们那里得到见证，才这样写。因为大约130年前，\Place{大罗马}和我们城市的主教们，曾写信谴责那些说子是受造物、不与父\OriginalTerm{同性同体}{coessential}的人。\Person{优西比乌}知道这事，他原是\Place{凯撒勒亚}的主教，起初是亚略异端的同伙；但后来，在\Place{尼西亚}公会议上签了字，写信给他自己的人，确认如下：\Quote{我们知道某些有口才且杰出的主教和作家，甚至在古代，就在论及父与子的天主性时使用了\OriginalTerm{同性同体}{coessential}一词。}\par

\Emph{主张子乃受造物，乃既不一致又站不住脚。}\par

那么，他们为何一再引证他们曾被罢黜的\Place{阿里米尼}会议？又为何拒绝\Place{尼西亚}会议，在那次会议上，他们的教父们签署了宣认，声明圣子属于圣父的性体，并与父\OriginalTerm{同性同体}{coessential}？他们为何四处奔走？因为如今，他们不仅与在尼西亚集会的众主教为敌，也与他们自己的伟大主教们和他们的朋友们为敌。那么，他们是谁的后嗣或继承者？他们怎能称那些人为父亲，却不接受他们妥善而宗徒性地拟定的宣认？因为，如果他们认为自己可以反对它，就让他们说出来，或者毋宁说作出回答，好叫他们被定罪，显明他们自相矛盾，当他们被问及是否相信圣子所说\BibleQuote{我与父原是一体}，\BibleQuote{看见我的，就看见了父}时。他们必然回答：\Quote{是的，经上记载，我们就信了。}但如果他们被问及这是怎样的一体，以及看见子的如何就是看见了父，当然，我们猜想他们会说：\Quote{由于相似}，除非他们已完全认同那些持与他们相似观点、被称为\OriginalTerm{阿诺摩派}{Anomœans}的人。但是，如果再问他们：\Quote{祂是怎样相似的？}他们就会厚颜无耻地说：\Quote{藉著完美的德行与和谐，藉著与父有同一旨意，不愿父所不愿的。}但是，要让他们明白，一个凭借德行和意志与天主相似的人，也可能改变主意；但圣言并非如此，除非祂只在部分上\Quote{相似}，且如同我们一般，因为祂在性体上也不像[天主]。但这些特征属于我们这些有起源、具受造本性的人。因为我们自己，虽然不能在性体上成为与天主相似，却可在德行上进步，效法天主，主赐予我们这恩宠，正如经上所说：\BibleQuote{你们应当慈悲，就像你们的父那样慈悲}：\BibleQuote{你们应当是成全的，如同你们的天父是成全的一样}\BibleRef{路 6:36}。至于受造物是可变的，无人能否认，因为天使犯了罪，亚当违了命，而一切受造都需要圣言的恩宠。但是，可变的事物无法与真正不变的天主相似，就如受造物不能与它的创造者相似一样。这就是为什么，关于我们，圣者说：\BibleQuote{主，谁能相似你}，以及\BibleQuote{主，众神之中谁能像你}；他所说的神，是指那些虽然受造，却已分受了圣言的人，正如祂自己说过：\BibleQuote{如果祂称那些承受天主话的人为神}\BibleRef{若 10:35}。但是，分受者不可能与所分受的那位相同或相似。例如，祂论到自己说：\BibleQuote{我与父原是一体}，暗示受造物并非如此。我们要问那些援引阿里米尼会议的人，一个受造的性体能否说：\BibleQuote{我见父所作的，我也照样作}。因为受造物是被造的，而不能创造；否则它们甚至创造了自己。如果，如他们所说，子是受造物，父是祂的创造者，那么子就必定是祂自己的创造者，因为祂能作父所做的事，正如祂所说的。但是，这种假设是荒谬且完全站不住脚的，因为无人能创造自己。\par

8. \Emph{圣子与圣父的关系是本质的，而非仅是伦理的。}\par

再请问他们，受造物能否说：\BibleQuote{凡父所有的一切，都是我的}\BibleRef{若 16:15}。如今，祂拥有创造和制造的专权、永恒、全能、不变性的专权。但是受造物不能有制造的能力，因为它们是受造物；也不能有永恒，因为它们的存在有开端；也不能有全能与不变性，因为它们受支配，且本性可变，如圣经所说。那么，如果这些专权属于圣子，显然不是因为祂的德行，如上面所说，而是本质上的，正如大公会议所说：\Quote{祂并非来自另一个性体}，而是属于父的，这些专权乃父所固有。然而，除了\OriginalTerm{同性同体}{coessential}外，那专属于父的性体、且是从父性体所出之子嗣的，还能是什么，或者说，我们还能赋予它别的名称吗？因为人在父内所见的，在子内也同样看见；并且不是由于分受，而是由于本质。这就是\BibleQuote{我与父原是一体}，以及\BibleQuote{看见我的，就看见了父}的意义。在这里，特别再次容易地显明他们的愚昧。如果你们认为子是凭借德行、意愿与不意愿的先决条件以及道德上的进步而与父相似，而这一切都归入性质的范畴；那么你们显然称天主为由性质与性体组合而成的。但谁还会容忍你们如此说呢？因为赋予一切存在、组合万物的天主，本身并非组合体，其本性也与祂藉圣言所造的万物不相似。千万不可有此想法。因为祂是单纯的性体，其中不存在性质，也没有，如雅各伯所说：\BibleQuote{任何变化或转动的阴影}\BibleRef{雅 1:17}。因此，如果证明了不是出于德行（因为在天主内没有性质，在子内也没有），那么圣子必定专属于天主的性体。只要你们的心智尚未彻底毁坏，你们一定会承认这一点。但是，那专属于并等同于天主性体，且由父本质所生的子嗣，若不是正因为此事实而与生祂的那位\OriginalTerm{同性同体}{coessential}，又是什么呢？因为这就是子对父的独特关系，凡否认此点者，就不认为圣言在本质与真理上是子。\par

9. \Emph{对\OriginalTerm{亚略主义}{Arianism}的诚实拒斥包含对\OriginalTerm{尼西亚准则}{Nicene test}的接受。}\par

教父们正是察觉到这一点，所以写下圣子与父\OriginalTerm{同质}{coessential}，并绝罚那些说\Quote{圣子有不同\OriginalTerm{实体}{Subsistence}}的人；他们并非自己发明言辞，而是如我们所说，从先前的教父那里学来的。然而，经过上述证明，他们的\OriginalTerm{阿里米尼会议}{Ariminian Synod}是多余的，任何他们引用的有关信仰的其他会议也是多余的。因为\Place{尼西亞}会议就足够了，它与古代的主教也是一致的，他们的教父也在上面签了名，他们应当尊重，否则很难被视为基督徒。但是如果即使经过这些证明，经过古代主教的见证，以及他们自己教父的签名，他们却佯装无知，对\InnerQuote{同质}一词感到惊恐，那么就让他们用更简单的语言，真诚地说出并坚持圣子按本性就是子，并按会议所规定的，绝罚那些说天主子是受造物或所造之物，或是从无中而来，或者曾有一个时间祂不存在，以及说祂是可变的、能改变的、有另一个实体的人。这样，他们就能逃脱\OriginalTerm{亚略异端}{Arian heresy}。我们确信，当他们真诚地绝罚这些观点时，他们就\Emph{事实上}承认了圣子出自父的本体，并与父\OriginalTerm{同质}{coessential}。因为这就是为什么教父们在说了圣子是同质之后，立即又补充道：\Quote{但凡说祂是受造物，或是所造之物，或是从无中而来，或曾有一个时间祂不存在}，\OriginalTerm{天主教会}{Catholic Church}予以绝罚：正是为了通过这种方式，让人明白\InnerQuote{同质}一词所指的是这些事。而\InnerQuote{同质}的含义从圣子不是受造物或所造之物可知：因为说\InnerQuote{同质}的人并不认为\OriginalTerm{圣言}{Word}是受造物；绝罚上述观点的人，同时也就认信圣子与父\OriginalTerm{同质}{coessential}；称祂\InnerQuote{同质}的人，就是真正而真实地称祂为天主子；而称祂为真天主子的人，就理解了经文：\BibleQuote{我與父原是一體，} 和 \BibleQuote{看見我的，就是看見父。}\par

10. \Emph{本信的目的；对\Place{米兰}的\Person{奥克森提乌斯}提出警告}\par

本来更适宜详细写下这些。但既然我们是写信给你们这些知情人，我们就简明地口授了，祈愿和平的联系在众人中得到维护，并使所有在\OriginalTerm{天主教会}{Catholic Church}内的人能说同样的话、持守同样的事。我们并非意图教导，而是提醒你们。不仅我们自己写信，而且所有\Place{埃及}和\Place{利比亚}的主教，约有九十位，也写信。因为在这事上我们全都同心合意，并且若有人恰巧不在场，我们也常常互相代签。我们既怀有这样的心意，又恰巧正聚在一起，便写信给我们亲爱的\Place{大罗马}主教\Person{达玛稣}，报告那入侵\Place{米兰}教会的\Person{奥克森提乌斯}的事；即他不单参与了\OriginalTerm{亚略异端}{Arian heresy}，还被指控犯下许多罪行，这些罪行是他与那位同他一样不虔敬的\Person{额我略}一起犯下的；并且，我们一方面对迄今为止他尚未被罢免并逐出教会表示惊讶，另一方面，我们为[\Person{达玛稣}]的虔诚以及那些在\Place{大罗马}聚集之人的虔诚而感谢，因为他们驱逐了\Person{乌尔撒奇}、\Person{瓦伦斯}以及与他们同伙的人，从而维护了\OriginalTerm{天主教会}{Catholic Church}的共融。我们祈求这共融在你们中间也得保全，因此恳求你们，如上所述，不要容忍那些提出一连串关于信仰的会议的人，那些会议分别在\Place{阿里米尼}、\Place{西尔米翁}、\Place{依索里亚}、\Place{色雷斯}、\Place{君士坦丁堡}举行，以及在\Place{安提约基}举行的许多不合规的会议。唯愿教父们在\Place{尼西亞}所宣认的信仰在你们中间独行其道，那次会议所有的教父，包括如今反对这信仰之人的教父，如上所述，都在场并签了名：这样，宗徒论到我们时也可以说：\BibleQuote{如今我称讚你们，因为你们在一切事上都记念我，并照我所传授给你们的，持守那些传统}（\BibleRef{格前 11:2}）。\par

11. \Emph{\OriginalTerm{圣神}{Holy Spirit}的\OriginalTerm{天主性}{Godhead}也包含在\OriginalTerm{尼西亞信经}{Nicene Creed}中}\par

因为这\Place{尼西亞}会议确实是对一切异端的禁令。它也推翻了那些亵渎\OriginalTerm{圣神}{Holy Spirit}、称祂为\OriginalTerm{受造物}{Creature}的人。因为教父们在讲论对圣子的信仰之后，立即加上了\Quote{我们信圣神}，为的是通过完美且完全地宣认对\OriginalTerm{天主圣三}{Holy Trinity}的信仰，使基督信仰的明确形式以及\OriginalTerm{天主教会}{Catholic Church}的教导得以彰显。因为无论在你们中间还是在众人中间，这事都已显明，任何基督徒都不可对此心存疑虑，即我们的信仰不在于受造物，而在于唯一的天主，全能的父，一切可见与不可见之物的创造者；在于唯一的主耶稣基督，祂的独生子；在于唯一的圣神；这是一位天主，在圣而完美的天主圣三中被认识，我们受洗进入其中，并在其中与\OriginalTerm{天主性}{Deity}结合，我们相信，在我们主基督耶稣内，我们也已继承了天国，藉着祂，愿光荣与权能归于父，至于无穷之世。阿们。\par

\SourceRecord{Translated by Archibald Robertson.；Nicene and Post-Nicene Fathers, Second Series；Vol. 4.；Edited by Philip Schaff and Henry Wace.；Buffalo, NY: Christian Literature Publishing Co.,；1892.；<http://www.newadvent.org/fathers/2819.htm>.}
